% !TeX encoding = UTF-8
% El diseño está separado en hvtbook.cls. Aquí solo editas el contenido.
% ARCHIVO PRINCIPAL: abra y compile este archivo en TeXstudio.
% Los estilos visuales (incluida la tabla actualizada) están en hvtbook.cls.
\documentclass{hvtbook}

\HVTTitulo{Prueba editable en LaTeX}
\HVTSubtitulo{Una muestra para comprobar si editar el cuadernillo resulta cómodo y si la web conserva su diseño.}
\HVTNivel{Nivel de prueba}
\HVTNumero{00}
\HVTPortada{../images/Microbit.png}{Imagen de portada utilizada únicamente en esta prueba.}

\begin{document}
\HVTMakeTitle

\begin{HVTPage}{Capítulo 1}{Editar sin tocar HTML}{Objetivo de la prueba}

\HVTLead{Al cambiar estas frases en TeXstudio y volver a generar, se actualizan tanto el PDF como la página web.}

\begin{HVTGrid}
  \begin{HVTColumn}{Lo que vamos a comprobar}
    \begin{itemize}
      \item Que las tildes y la letra ñ funcionen correctamente.
      \item Que las listas sean fáciles de reorganizar.
      \item Que el diseño web conserve sus hojas y colores.
    \end{itemize}
  \end{HVTColumn}

  \begin{HVTColumn}{Cambios sugeridos}
    \begin{itemize}
      \item Cambia el título de esta columna.
      \item Agrega un elemento nuevo a esta lista.
      \item Modifica el texto de la siguiente figura.
    \end{itemize}
  \end{HVTColumn}
\end{HVTGrid}

\HVTImagen{../images/Microbit.png}{Figura de prueba. Este pie se edita directamente en LaTeX.}

\end{HVTPage}

\begin{HVTPage}{Práctica}{Haz un cambio pequeño}{Prueba la edición en TeXstudio}

\begin{HVTExample}{Ejemplo 1}{Mi primer cambio}
Abre este archivo en TeXstudio, cambia una frase y guarda. Después ejecuta el generador para verla en la web y en el PDF.
\end{HVTExample}

\begin{HVTSteps}
  \item \textbf{Cambia el título.} Modifica el contenido de \texttt{\textbackslash HVTTitulo}.
  \item \textbf{Edita una lista.} Agrega o elimina cualquier elemento escrito con \texttt{\textbackslash item}.
  \item \textbf{Genera los resultados.} Ejecuta \texttt{generar-prueba.cmd}.
  \item \textbf{Comprueba el diseño.} Abre \texttt{cuadernillo-prueba-web.pdf}.
\end{HVTSteps}

\end{HVTPage}

\begin{HVTPage}{Componentes}{Fórmulas, diagramas y figuras}{Balance de un sistema renovable}

\HVTLead{Esta hoja permite aprobar cómo se presentan relaciones matemáticas, procesos técnicos e imágenes antes de migrar los cuadernillos reales.}

\begin{HVTFlow}
  \HVTFlowItem{1}{Recurso}{Viento, agua o radiación disponible.}
  \HVTFlowItem{2}{Conversión}{El equipo transforma el recurso en energía útil.}
  \HVTFlowItem{3}{Servicio}{La energía cubre una demanda identificada.}
\end{HVTFlow}

\HVTEquation{\frac{1}{4} - \frac{x-3}{16} + \frac{1}{64}(x-3)^2 - \frac{1}{256}(x-3)^3 + \frac{(x-3)^4}{1024} - \frac{(x-3)^5}{4096} + \mathcal{O}((x-3)^6)}{Ejemplo de una expansión presentada como ecuación independiente, con fracciones apiladas y potencias.}

\begin{HVTGallery}
  \HVTImagen{../images/Esp32.png}{Figura A. Controlador usado para registrar las variables del sistema.}
  \HVTImagen{../images/Electrolizador.png}{Figura B. Equipo técnico representado dentro del mismo lenguaje visual.}
\end{HVTGallery}

\end{HVTPage}

\begin{HVTPage}{Validación}{Datos, seguridad y fuentes}{Ficha técnica de una medición}

\begin{HVTTable}{Medición}{Viento}{Agua}
  \HVTTableRow{Entrada}{Velocidad, dirección, densidad y turbulencia.}{Caudal, niveles, altura neta y sedimentos.}
  \HVTTableRow{Mecánica}{rpm, par, vibración, deformación y temperatura.}{rpm, par, vibración, presión y temperatura.}
  \HVTTableRow{Salida}{Potencia, energía, tensión, corriente o servicio mecánico.}{Potencia, energía, tensión, corriente o servicio mecánico.}
  \HVTTableRow{Estado}{Humedad de madera, fisuras, uniones y recubrimiento.}{Desgaste, corrosión, obstrucción, madera y apoyos.}
\end{HVTTable}

\begin{HVTWarning}{Puerta de decisión}
No seleccionar equipos antes de contar con mediciones representativas, revisión de seguridad y condiciones de mantenimiento.
\end{HVTWarning}

\end{HVTPage}

\begin{HVTPage}{Evidencia}{Actividad, código y fuentes}{Registro reproducible}

\begin{HVTTask}{Actividad de comprobación}
  \item Registrar fecha, lugar, instrumento y responsable.
  \item Anotar unidades, incertidumbre y condiciones ambientales.
  \item Comparar el resultado con la demanda identificada.
\end{HVTTask}

\begin{HVTCode}{Pseudocódigo de validación}
leer medicion
si medicion esta en rango:
    guardar dato
si no:
    registrar alerta
\end{HVTCode}

\begin{HVTReferences}
  \HVTReference{Guía de instrumentos y métodos de observación de la OMM}{https://community.wmo.int/}
  \HVTReference{Material educativo de la Fundación}{https://hidrogenoverdeturquesa.com/fundacion}
\end{HVTReferences}

\end{HVTPage}

\end{document}
